\documentclass[11pt]{article}

\usepackage[margin=1in]{geometry}
\usepackage{amsmath,amssymb,amsthm}
\usepackage[hidelinks]{hyperref}
\hypersetup{
  pdftitle={Aligned-Center Criteria for Products with Multiple q-Integer Spacers},
  pdfauthor={Hainan Zhao}
}

\newtheorem{theorem}{Theorem}
\newtheorem{lemma}[theorem]{Lemma}
\newcommand{\qint}[1]{[#1]_q}
\newcommand{\qrint}[2]{[#1]_{q^{#2}}}

\title{Aligned-Center Criteria for Products with Multiple
\(q\)-Integer Spacers}
\author{Hainan Zhao}
\date{August 2026}

\begin{document}
\maketitle

\begin{abstract}
Connelly, Ito, Martinez, Shevchenko, and Yang conjectured a sufficient
condition for the unimodality of a product of ordinary \(q\)-integers and
one lacunary factor.  We first prove that condition for every number of
ordinary factors and every spacer step in its stated range.  We then prove
a hybrid criterion for arbitrarily many lacunary factors: disjoint
divisibility pairs are absorbed factor by factor, and the remaining spacer
increments are assigned to the remaining ordinary factors by a nonnegative
integer allocation matrix.  No coprimality among the spacer steps is
required.  Both results follow from one aligned-center recursion.  The
criterion is sufficient rather than necessary, and explicit small examples
map that boundary.
\end{abstract}

\section{Introduction}

For a positive integer \(n\), write
\[
 [n]_q=1+q+\cdots+q^{n-1}.
\]
Connelly, Ito, Martinez, Shevchenko, and Yang proposed the following
criterion \cite[Conjecture~5.4]{CIMSY}.  If \(r\geq2\), \(k\geq1\), and
\(a_1,\ldots,a_k,b\) are positive integers satisfying
\[
 r\mid a_i\quad\text{for some }i,
 \qquad\text{or}\qquad
 b\leq1+\sum_{i=1}^k\left\lfloor\frac{a_i}{r}\right\rfloor,
 \tag{1}
\]
then
\[
 \prod_{i=1}^k[a_i]_q[b]_{q^r}                    \tag{2}
\]
is unimodal.  They conjectured necessity when \(k\leq3\) or \(r\leq3\),
having already proved the full equivalence for \(k=1\) and for
\((k,r)=(2,2)\).  Their bounded checks are separate computational evidence.

Our first result settles the sufficient direction exactly as stated.

\begin{theorem}[one-spacer criterion]\label{thm:one-spacer}
Let \(r\geq2\), \(k\geq1\), and let \(a_1,\ldots,a_k,b\) be positive
integers.  If either condition in \emph{(1)} holds, then the polynomial
\emph{(2)} is symmetric and unimodal.
\end{theorem}

The same mechanism permits several spacers.  For the next statement, choose
sets \(J\subseteq\{1,\ldots,s\}\) and
\(I\subseteq\{1,\ldots,k\}\), and a bijection
\(\phi:J\to I\).  Write \(J^c\) and \(I^c\) for their complements.

\begin{theorem}[hybrid multi-spacer criterion]\label{thm:hybrid}
Let \(k,s\geq1\); let \(a_1,\ldots,a_k,b_1,\ldots,b_s\) be positive
integers; and let \(r_1,\ldots,r_s\geq1\).  Suppose that
\[
 r_j\mid a_{\phi(j)}\qquad(j\in J),
\]
and that there are nonnegative integers \(d_{j,i}\), for
\(j\in J^c\) and \(i\in I^c\), such that
\[
 \sum_{i\in I^c}d_{j,i}=b_j-1\qquad(j\in J^c),       \tag{3}
\]
and
\[
 a_i-\sum_{j\in J^c}r_jd_{j,i}\geq1
 \qquad(i\in I^c).                                  \tag{4}
\]
Then
\[
 \prod_{i=1}^k[a_i]_q\prod_{j=1}^s[b_j]_{q^{r_j}}
\]
is symmetric and unimodal.  Empty residual index sets are allowed, with
the displayed conditions interpreted as empty sums.
\end{theorem}

Theorem~\ref{thm:hybrid} has no coprimality hypothesis.  It is genuinely
beyond the named conjecture because it allows arbitrarily many spacers, but
it remains only sufficient.  The source example
\(([3]_q)^4[2]_{q^4}\) is unimodal while violating (1), so no general
necessity claim is made.  At two spacers,
\([2]_{q^2}[2]_{q^3}=1+q^2+q^3+q^5\) is symmetric but not unimodal and is
excluded by (4).  The nearby unimodal product
\([5]_q[2]_{q^2}[2]_{q^3}\) is also outside (4), showing that the new
condition is not necessary.  We do not give an exact characterization of
multi-spacer unimodality.

\section{Symmetric unimodality}

A finite nonnegative coefficient sequence is \emph{symmetric unimodal about
\(C\)} if coefficients at exponents equidistant from \(C\) agree and the
coefficients weakly increase up to \(C\).  We allow initial and terminal
zeros; this convention makes translates convenient.

\begin{lemma}\label{lem:closure}
If \(F\) and \(G\) are nonnegative symmetric unimodal polynomials about
centers \(C_1\) and \(C_2\), respectively, then \(FG\) is symmetric
unimodal about \(C_1+C_2\).  A sum of nonnegative symmetric unimodal
polynomials having a common center \(C\) is symmetric unimodal about \(C\).
\end{lemma}

\begin{proof}
The assertion about sums follows coefficientwise.  For products, decompose
each symmetric unimodal coefficient sequence into a nonnegative linear
combination of indicator sequences of nested intervals having the same
center.  Unimodality permits zeros only before and after the nonzero support,
not internally, which licenses this nested-interval decomposition.  The
convolution of two interval indicators is a symmetric
trapezoidal sequence, hence unimodal, and all resulting convolutions have
the sum of the original centers.  Summing them proves the claim.
\end{proof}

\section{The aligned-center recursion}

\begin{lemma}[aligned-center recursion]\label{lem:recursion}
For positive integers \(a,b,r\),
\[
 [a+r]_q[b+1]_{q^r}
 =q^r[a]_q[b]_{q^r}+[a+r(b+1)]_q.                  \tag{5}
\]
\end{lemma}

\begin{proof}
Use
\[
 [a+r]_q=[r]_q+q^r[a]_q,
 \qquad
 [b+1]_{q^r}=1+q^r[b]_{q^r}.
\]
After subtracting the first term on the right side of (5), the remaining
terms are
\begin{align*}
 [r]_q[b+1]_{q^r}+q^r[a]_q
       \bigl(1+(q^r-1)[b]_{q^r}\bigr)
 &= [r(b+1)]_q+q^{r(b+1)}[a]_q\\
 &= [a+r(b+1)]_q.
\end{align*}

For a coefficient-level proof independent of this manipulation, the left
side counts pairs \((x,j)\) with \(0\leq x<a+r\) and \(0\leq j\leq b\),
weighted by \(x+rj\).  The pairs with \(x<a\) and \(j\geq1\) give
\(q^r[a]_q[b]_{q^r}\).  The complementary pairs have either \(j=0\), or
\(x\geq a\) and \(j\geq1\); their weights occur once each and fill the
consecutive interval from \(0\) to \(a+r(b+1)-1\).  They give the second
term in (5).
\end{proof}

\section{The hybrid criterion}

We first prove the allocation part of Theorem~\ref{thm:hybrid}.

\begin{lemma}[matrix induction]\label{lem:matrix}
Let \(k\geq1\) and \(s\geq0\).  If nonnegative integers \(d_{j,i}\)
satisfy (3)--(4) with \(I=J=\varnothing\), then
\[
 \prod_{i=1}^k[a_i]_q\prod_{j=1}^s[b_j]_{q^{r_j}}
\]
is symmetric and unimodal.
\end{lemma}

\begin{proof}
Induct on \(s\).  The case \(s=0\) is a product of ordinary
\(q\)-integers and follows from Lemma~\ref{lem:closure}.  For \(s\geq1\),
put
\[
 a_i^{(0)}=a_i-r_1d_{1,i}.
\]
Condition (4) shows that the rows \(2,\ldots,s\) certify the induction
hypothesis for
\[
 \prod_i[a_i^{(0)}]_q\prod_{j=2}^s[b_j]_{q^{r_j}}.
\]
This is the base at first-spacer length one.

Reconstruct the \(b_1-1\) increments in any order realizing the first row
of the matrix.  At one step, let \(x\) be the current selected ordinary
length, \(c\) the current first-spacer length, \(B(q)\) the product of the
other ordinary factors, and \(R(q)\) the product of spacers
\(2,\ldots,s\).  Lemma~\ref{lem:recursion} gives
\begin{align*}
 &[x+r_1]_qB(q)[c+1]_{q^{r_1}}R(q)\\
 &\qquad=q^{r_1}[x]_qB(q)[c]_{q^{r_1}}R(q)
 +[x+r_1(c+1)]_qB(q)R(q).                         \tag{6}
\end{align*}
The first summand is a translate of the preceding polynomial.  In the
second, every ordinary length is at least its base value, and the selected
length is larger still.  The unchanged rows \(2,\ldots,s\) therefore
continue to satisfy (4), so the induction hypothesis makes the correction
term symmetric unimodal.

If the old degree is
\[
 E=(x-1)+\deg B+r_1(c-1)+\deg R,
\]
the new degree is \(E+2r_1\).  The translated old support has endpoints
\(r_1\) and \(E+r_1\), while the correction term has degree
\[
 x+r_1(c+1)-1+\deg B+\deg R=E+2r_1.
\]
Both summands in (6) therefore have center \((E+2r_1)/2\).  Their sum is
symmetric unimodal by Lemma~\ref{lem:closure}.  Iteration reconstructs the
first spacer and all ordinary lengths, completing the induction.
\end{proof}

\begin{proof}[Proof of Theorem~\ref{thm:hybrid}]
For every \(j\in J\), write \(a_{\phi(j)}=r_je_j\) and set \(z=q^{r_j}\).
Then
\[
 [a_{\phi(j)}]_q[b_j]_{q^{r_j}}
 =[r_j]_q\bigl([e_j]_z[b_j]_z\bigr)\big|_{z=q^{r_j}}. \tag{7}
\]
The product \([e_j]_z[b_j]_z\) is symmetric unimodal in \(z\) by
Lemma~\ref{lem:closure}.  Substitution followed by multiplication by
\([r_j]_q\) repeats each coefficient in a consecutive block of length
\(r_j\), so (7) is symmetric unimodal in \(q\).

The map \(\phi\) is a bijection between the selected sets, so each ordinary
factor is consumed by at most one spacer.  Thus the factors in (7) are
disjoint and do not interact with the residual allocation matrix.  The
residual product is symmetric unimodal by Lemma~\ref{lem:matrix}; multiplying
it by the pair factors preserves symmetric unimodality by
Lemma~\ref{lem:closure}.

If no residual ordinary factors remain, (3) forces every residual spacer
length to equal one, so that residual product is the constant polynomial.
This also covers the empty case.
\end{proof}

The statement permits \(r_j=1\).  Such a spacer is an ordinary bracket.  In
the one-factor, one-spacer matrix case, (3)--(4) become \(b\leq a\); swapping
the two ordinary brackets covers the opposite ordering.  No step of the
proof assumes \(r_j\geq2\).

\section{The named conjecture as a corollary}

\begin{proof}[Proof of Theorem~\ref{thm:one-spacer}]
If \(r\mid a_i\) for some \(i\), pair that ordinary factor with the sole
spacer in Theorem~\ref{thm:hybrid}.

Otherwise, put \(c_i=\lfloor a_i/r\rfloor\).  The inequality in (1) permits
integers \(d_i\) with
\[
 0\leq d_i\leq c_i,
 \qquad \sum_i d_i=b-1.
\]
Because \(r\nmid a_i\),
\[
 a_i-rd_i\geq a_i-r\left\lfloor\frac{a_i}{r}\right\rfloor
 =a_i\bmod r\geq1.                                \tag{8}
\]
Thus the single row \((d_i)_i\) satisfies (3)--(4), and
Theorem~\ref{thm:hybrid} applies.
\end{proof}

\section{Sharpness and relation to width five}

The allocation condition pays for every spacer increment against the
worst-case remaining capacity.  For
\([a]_q[2]_{q^2}[2]_{q^3}\), the unique one-column allocation requires
\(a-2-3\geq1\), so it activates at \(a=6\).  The exact products at
\(a=4\) and \(a=5\) have coefficient sequences
\[
 (1,1,2,3,2,3,2,1,1),\qquad
 (1,1,2,3,3,3,3,2,1,1),
\]
respectively.  Hence \(a=4\) is not unimodal, while \(a=5\) is unimodal
and violates (4); a mild adaptive refinement certifies it but is omitted
because it adds no width-five class.  For repeated steps
\([a]_q[2]_{q^2}[2]_{q^2}\), the allocation threshold is \(a=5\), while
\(a=2\) already has coefficients \((1,1,2,2,1,1)\).  These examples show
why no necessity or sharp-threshold conjecture is asserted.

There is nevertheless a stable bounded pattern in the one-spacer necessity
regimes.  For \(k\leq5\), \(r\leq6\), and all parameters at most 15, an
exact census of the 33,728 rows violating (1) with \(k\leq3\) or \(r\leq3\)
found a midpoint-side dip in every row.  Across all 4,576 fixed
\((r,a_1,\ldots,a_k)\) groups, the location of the first dip was independent
of \(b\) throughout the violating range, and its depth was always one or
two.  This is computational evidence, not an exact-boundary conjecture.

The width-five \(q\)-Fibonomial theorem requires three spacer steps
\(2,3,5\) in the classes admitting a bracket decomposition.  Across the 60
stable residue classes, matrix allocation reaches 18 classes, divisibility
absorption alone reaches none, and the hybrid reaches the same 18.  It
misses 24 further decomposable classes and all 18 classes without such a
decomposition.  The equal counts are coincidental: the 18 reached classes
and the 18 undecomposable classes are disjoint.  Absorption alone cannot
cover a stable class because all three spacers are nontrivial while only two
ordinary factors remain; allowing absorption inside the hybrid adds no new
class.  Its role is nevertheless essential in
Theorem~\ref{thm:one-spacer}, where it recovers the divisibility branch.
The same exact class census shows that the mild adaptive refinement above
does not enlarge the 18-class coverage.

Thus 24 of the 42 decomposable classes remain beyond the criterion.  What
additional mechanism reaches those 24 classes while retaining an
aligned-center proof?  A separate restricted-partition coefficient-
difference argument treats all width-five classes uniformly, but does not
answer that structural question.  An exact characterization of
multi-spacer unimodality also remains open.

\section{Reproducibility and scope}

The proofs use exact algebra.  A standard-library checker verifies (5) by
both polynomial multiplication and the independent pair partition, replays
the one-spacer induction, checks the matrix induction through four spacers,
tests divisibility absorption, and reproduces the sharpness and width-five
coverage calculations.  From the repository root run
\begin{verbatim}
python3 proof/qanalog_multispacer_criterion.py
\end{verbatim}
The finite corpus is regression evidence only; it is not an assumption in
either universal proof.

Theorems~\ref{thm:one-spacer} and \ref{thm:hybrid} are sufficiency results.
They do not characterize all unimodal products, prove necessity beyond
previously known cases, or settle general \(q\)-Fibonomial unimodality.

\paragraph{AI assistance.}
The author used AI assistance in developing the proof presentation and its
verification code.

\begin{thebibliography}{1}

\bibitem{CIMSY}
Brendan B. Connelly, Ezekiel Ito, Thomas C. Martinez, Olha Shevchenko, and
Kacey Yang.
\newblock Unimodality of \(q\)-Fibonomial coefficients for small cases.
\newblock \emph{arXiv:2605.12822v1 [math.CO]}, May 12, 2026.

\end{thebibliography}

\end{document}
